% Ad Atticum 14.22 --- headnote
% ad-atticum-14-22, 14 May 44 BC, in Puteolano

\textit{Cicero to Atticus, written at the Puteolan villa on 14
May 44 BC --- Perseus dateline \textit{Scr. in Puteolano prid.
Id. Mai. a. 710 (44)}. A brisk note dashed off because Pilia
has warned him that couriers are leaving for Rome on the Ides.
Cicero gives notice that he leaves Puteoli for Arpinum on 17
May, but wants first to ``scent out'' (\textit{odorari
diligentius}) what is coming. His own pupil --- that is,
Hirtius, due at dinner that very day --- is openly in love
with the man Brutus wounded, and the Caesarian camp's
working premise (\textit{[Greek: hupothesis]}) has been
laid bare: a great man was murdered, the state thrown into
confusion, his clemency was his ruin; everything will be
undone as soon as the Liberators stop being feared.}

\textit{Section 2 voices for the first time, in this
correspondence, what will become a recurring fear: that
Sextus Pompeius may come from Spain with a strong army
(\textit{[Greek: eulogon]}, ``reasonable'' enough to expect),
and that war will then certainly come. Cicero --- with three
densely packed Greek words within a single line: \textit{
[Greek: phainoprosop\=eteon]} (``one must put on a brave
face''), \textit{[Greek: iteon]} (``one must go''), and the
Sophoclean tag \textit{[Greek: allois en esthlois tond'
ap\=othountai psogon]} (a fragment of Sophocles' lost
\textit{Erigone}, applied to the Liberators) --- weighs and
rejects the camp option for himself. The Ides of March, he
admits, no longer console him as they did: they ``contain a
great flaw,'' namely Antony left alive. He may travel on his
votive embassy yet; he is being warned by many not to attend
the Senate on the Kalends of June, when troops are said to be
secretly mustered against the very men who will be safer
anywhere than there.}
